% !PRINTED 531
\weberpage{531}

\gesperrt{Whatever values one may assign to the many-valued algebraic quantities that occur in the solution, the expressions (15) thus always represent the roots of our biquadratic equation in some order.}

With regard to the remaining ways of solving the biquadratic equation we can be briefer.

If we first adjoin not $\sqrt{D}$ but a function belonging to the group $P_1$, we obtain a cubic resolvent which is not a normal equation. For this function we can choose, say, $y^2$, which satisfies the following cubic equation in that quantity:
\weq{16}{y^6 - 2Ay^4 + A^2y^2 - D = 0}
Better still, one takes as roots of the cubic resolvent
\weq{17}{z = y_1 - y_2, \quad z_1 = y_2 - y, \quad z_2 = y - y_1.}
Then $z$ belongs to the group $P_1$, as already follows from the second equation (10), or as may also be seen directly. From (8) and (9) one obtains for $z$ the cubic equation
\weq{18}{z^3 - 3Az + B = 0.}
If we further set
\weq{19}{u = \alpha\alpha_1 + \alpha_2\alpha_3}
then
\[
z = a_2 - 3u
\]
and hence obtain the resolvent for $u$:
\weq{20}{u^3 - a_2u^2 + (a_1a_3 - 4a_4)u + 4a_2a_4 - a_1^2a_4 - a_3^2 = 0.}

But if one adjoins not one but all the roots of equation (18) or (20), one arrives again at the group $Q_1$. The quantities $v_1, v_2, v_3$ are by (12) expressed immediately in terms of $z, z_1, z_2$, and one can therefore also apply the formulae (15) for the representation of the $\alpha$ by this procedure.

One can further arrive at the solution of the biquadratic equation by first adjoining a function of the roots belonging to the cyclic group $P_2$ (a cyclic function), as for example
\weq{21}{w = \alpha\alpha_1^2 + \alpha_1\alpha_2^2 + \alpha_2\alpha_3^2 + \alpha_3\alpha^2.}
This function is six-valued and is therefore the root of an equation of the sixth degree. But this equation of the sixth degree is, as one easily sees, imprimitive, and can be
%CONT>>
